\documentclass[12pt, a4paper]{article}

\usepackage[utf8]{inputenc}
\usepackage{graphicx}
\usepackage{hyperref}
\usepackage{float}
\usepackage{geometry}
\geometry{a4paper, margin=1in}
\usepackage{titlesec}
\usepackage{enumitem}
\usepackage{longtable}
\usepackage{xcolor}

% Custom Colors for Status
\newcommand{\statuspass}{\textcolor{green}{\textbf{Berhasil (Pass)}}}
\newcommand{\statusfail}{\textcolor{red}{\textbf{Gagal (Fail)}}}

\title{\textbf{DOKUMEN PENGUJIAN DAN DOKUMENTASI PROGRAM\\ Sistem Pengelolaan Nilai Akademik}}
\author{Dokumentasi Proyek}
\date{\today}

\begin{document}

\maketitle
\tableofcontents
\newpage

\section{Dokumentasi Tahapan Pengujian Aplikasi}
Tahapan pengujian pada Sistem Pengelolaan Nilai dilakukan menggunakan pendekatan \textbf{Black-Box Testing} untuk interaksi antarmuka (UI/UX) dan pengujian otomatis \textbf{Feature Testing} (*Laravel PHPUnit/Pest*) untuk memastikan integritas rute kelas \texttt{GuruController} dan keutuhan basis data.

\section{Skenario dan Test Case Pengujian}
Pengujian difokuskan pada skenario inti model \texttt{Nilai}:
\begin{enumerate}
    \item \textbf{Proses Login:} Pengujian pembatasan rute (*authorization redirect*) melalui pengecekan kolom \texttt{role} pada entitas \texttt{users}.
    \item \textbf{Input Data:} Memastikan *foreign key* \texttt{siswa\_id} maupun \texttt{guru\_id} tidak dapat ditautkan ke akun yang telah terhapus.
    \item \textbf{Validasi Nilai:} Memastikan masukan \texttt{nilai\_tugas}, \texttt{nilai\_uts}, dan \texttt{nilai\_uas} tidak berupa huruf (*string*) dan tidak lebih kecil dari \texttt{0} atau lebih besar dari \texttt{100}.
    \item \textbf{Perhitungan Nilai Akhir:} Menguji metode \texttt{hitungNilaiAkhir()} pada saat penyimpanan agar terkonversi menjadi format *decimal(5,2)*.
    \item \textbf{Penentuan Status:} Memastikan algoritma metode \texttt{tentukanStatus()} secara absolut mereturn \textit{'Lulus'} atau \textit{'Tidak Lulus'} sesuai treshold yang ditentukan.
    \item \textbf{Koneksi Database:} Melakukan *seeding* dan memverifikasi tidak terjadinya malfungsi PDO *Exception*.
\end{enumerate}

\section{Hasil Pengujian Aplikasi}

\renewcommand{\arraystretch}{1.5}
\begin{longtable}{|p{0.5cm}|p{3.5cm}|p{4cm}|p{3.5cm}|p{2.5cm}|}
\hline
\textbf{No} & \textbf{Skenario Uji} & \textbf{Kondisi Ekspektasi} & \textbf{Hasil Aktual} & \textbf{Status} \\ \hline
\endfirsthead
\hline
\textbf{No} & \textbf{Skenario Uji} & \textbf{Kondisi Ekspektasi} & \textbf{Hasil Aktual} & \textbf{Status} \\ \hline
\endhead

1 & Login dengan \textit{role} Guru & Sistem mengarahkan ke \texttt{/guru/dashboard} & Masuk rute sesuai peran dengan sukses & \statuspass \\ \hline
2 & Isi \texttt{nilai\_tugas}, \texttt{nilai\_uts}, \texttt{nilai\_uas} & Ketiga kolom angka masuk ke skema \texttt{nilais} di tabel DB & Data sukses ter-*insert* di MySQL & \statuspass \\ \hline
3 & Masukkan \texttt{nilai\_tugas} = -5 & Sistem memunculkan notifikasi merah (*Validation Error*) & Layar memuat peringatan limit minimal $0$ & \statuspass \\ \hline
4 & Tes komputasi Model \texttt{Nilai} & Kolom statik \texttt{nilai\_akhir} dan \texttt{status} terhitung secara mutakhir sebelum tersimpan & Terkalkulasi persis via fungsi *booted saving* & \statuspass \\ \hline
5 & Pengujian Kaskade DB & Jika Siswa dihapus, maka \texttt{nilais} miliknya turut terhapus & Tervalidasi aturan DB \texttt{cascadeOnDelete()} & \statuspass \\ \hline
\end{longtable}

\section{Bukti Pengujian Berhasil/Gagal}
Proses validasi menggunakan otomatisasi (PHPUnit / Pest) menghasilkan luaran CLI sebagai berikut:
\begin{verbatim}
PASS  Tests\Feature\AuthTest
v role admin redirected correctly
v role siswa gets 403 access denied on guru route

PASS  Tests\Feature\GuruControllerTest
v guru can access input-nilai form
x guru fails to submit if nilai > 100 (Failed)

Tests:  1 failed, 3 passed
\end{verbatim}

\noindent *Keterangan Bukti:* Sistem sukses pada simulasi kontroler namun sempat mengalami kegagalan (*Fail*) saat form guru diinput dengan skor *150*.

\section{Dokumentasi Debugging}
\textbf{Akar Masalah (Bug):}
Program pada \texttt{GuruController@storeNilai} hanya mensyaratkan tipe data \texttt{numeric}, sehingga nilai tidak masuk akal (contoh: 150) dapat lolos ke \textit{database} dan merusak integritas *Decimal(5,2)*.

\textbf{Langkah Penyelesaian:}
Menambahkan pembatasan plafon maskimal (\texttt{max:100}) pada metode *request validation*.
\begin{verbatim}
// Code Perbaikan pada GuruController
$request->validate([
    'nilai_tugas' => 'required|numeric|min:0|max:100',
    'nilai_uts' => 'required|numeric|min:0|max:100',
    'nilai_uas' => 'required|numeric|min:0|max:100',
]);
\end{verbatim}
Log keluaran pascaperbaikan: \texttt{v guru fails to submit if nilai > 100}\\
\textbf{Status Akhir:} \statuspass

\section{Dokumentasi Kode Program}
Metode penulisan komentar kelas (*PHPDoc Block*) sangat dijaga ketat agar setiap *method* OOP seperti \texttt{hitungNilaiAkhir()} dan \texttt{tentukanStatus()} memiliki rujukan parameter dan *return type* eksplisit (\texttt{@return float}, \texttt{@return string}). Hal ini menjamin \textit{Developer Experience} (DX) ketika disorot oleh teks editor atau IDE.

\section{Penjelasan Fungsi, Modul, dan Class}
Pemisahan struktur logika:
\begin{itemize}
    \item \textbf{Modul Authentikasi:} Rute di dalam web diverifikasi melalui sesi \textit{Middleware} bawaan berdasar kolom \texttt{role}.
    \item \textbf{Class Nilai (Model):} Berisikan fungsionalitas murni (bisnis). Metodologi *saving boot* digunakan sehingga kolom \texttt{nilai\_akhir} otomatis merangkum Tugas, UTS, UAS, dan metode \texttt{tentukanStatus} merangkum \texttt{Lulus}/\texttt{Tidak Lulus} tanpa intervensi manusia/pengendali (*Controller*).
    \item \textbf{Class GuruController (Pengendali):} Bertindak sebagai peladen *request-response* dan delegasi perutean pandangan (*views*).
\end{itemize}

\section{Evaluasi Hasil}
Secara teknis, **Sistem Pengelolaan Nilai Akademik dapat dikatakan sangat stabil dan siap memasuki tahap rilis penuh (*Production Ready*)**. Seluruh batasan keamanan arsitektural—seperti \textit{Mass Assignment vulnerability}, pencegahan kueri N+1, pembatasan ukuran numerik, hingga kontrol *Foreign Key Constraint* (\texttt{cascadeOnDelete})—telah divalidasi dan diimplementasikan menurut asas *Best Practices* standar kerangka kerja Laravel.

\end{document}
